\chapter{Primati: i nostri cugini troppo simili}

Gli scimpanzé condividono con noi il 98.8\% del DNA, e si vede. Nei test di laboratorio, mostrano molti dei nostri stessi bias cognitivi: avversione alle perdite, effetto ancoraggio, bias di conferma. È come guardare noi stessi allo specchio, solo con più peli e meno pretese di razionalità.

Ma c'è una differenza cruciale. Uno scimpanzé che ha paura irrazionale di un certo tipo di frutto al massimo avrà una dieta un po' limitata. Un umano che ha paura irrazionale degli OGM può influenzare le politiche alimentari di intere nazioni.

Il nostro problema non è avere un cervello di scimmia. Il problema è avere un cervello di scimmia con il potere di costruire armi nucleari, manipolare l'economia globale, e influenzare milioni di persone attraverso i social media.

Abbiamo la neurobiologia di una specie che viveva in gruppi di 50 individui, ma il potere di una specie che domina un pianeta. È un po' come dare le chiavi di una Ferrari a qualcuno che ha imparato a guidare sul trattore del nonno.

\bigskip

Ma andiamo con ordine, e vediamo quanto siamo davvero simili ai nostri cugini pelosi. Gli esperimenti sui primati degli ultimi decenni hanno rivelato somiglianze così inquietanti che a volte ci chiediamo se non stiamo semplicemente studiando versioni più semplici di noi stessi.

Prendiamo l'effetto ancoraggio, quel bias che fa sì che il primo numero che sentiamo influenzi tutte le nostre valutazioni successive. Negli esperimenti, agli scimpanzé viene mostrato un numero casuale, poi viene chiesto loro di scegliere tra diverse quantità di cibo. Incredibilmente, le loro scelte sono sistematicamente influenzate dal numero che hanno visto per primo, esattamente come succede agli esseri umani.

È come se il cervello dei primati - il nostro incluso - fosse programmato per usare qualsiasi informazione disponibile come punto di riferimento, anche quando quella informazione è completamente irrilevante. È un meccanismo che probabilmente serviva per fare stime rapide nell'ambiente ancestrale: se vedi tre frutti su un albero, ti aspetti di trovarne un numero simile su alberi dello stesso tipo.

Il problema è che questo meccanismo si attiva anche quando non dovrebbe. Lo scimpanzé che ha visto il numero 8 tenderà a scegliere quantità di cibo più vicine a 8, anche se il numero non ha nulla a che fare con il cibo. L'essere umano che ha sentito parlare di un prezzo di 1000 euro per un prodotto tenderà a considerare "ragionevole" un prezzo di 600 euro, anche se il prodotto vale molto meno.

È lo stesso bias, la stessa logica cognitiva, le stesse conseguenze irrazionali. La differenza è che le conseguenze dello scimpanzé rimangono confinate al laboratorio, mentre le nostre si estendono all'economia globale.

\medskip

Ancora più impressionante è come i primati gestiscono le perdite. Negli esperimenti di avversione alle perdite, ai bonobo vengono offerte due scelte: una ricompensa sicura di due frutti, o una ricompensa incerta che può essere di quattro frutti o zero frutti. Il valore atteso della seconda opzione è migliore (due frutti in media), ma i bonobo preferiscono sistematicamente la certezza di due frutti.

È esattamente il comportamento umano: preferiamo la certezza di un piccolo guadagno all'incertezza di un guadagno potenzialmente maggiore. Ma quando lo stesso esperimento viene fatto in termini di perdite - perdere sicuramente due frutti o rischiare di perderne quattro o zero - improvvisamente i bonobo diventano giocatori d'azzardo, proprio come noi.

Questo non è apprendimento o cultura: è neurobiologia pura. È il modo in cui il cervello dei primati - tutti i primati, noi inclusi - valuta rischi e ricompense. È stato programmato così da milioni di anni di evoluzione, e funziona secondo la stessa logica sia nel cervello di un bonobo che in quello di un trader di Wall Street.

La differenza è che quando un bonobo prende una decisione irrazionale sulle perdite, al massimo perde qualche frutto. Quando un trader umano prende la stessa decisione irrazionale, può mandare in rovina milioni di pensionati.

\bigskip

Ma forse l'aspetto più affascinante - e preoccupante - della nostra somiglianza con gli altri primati riguarda i bias sociali. Gli scimpanzé mostrano chiari segni di pregiudizio verso il gruppo di appartenenza: preferiscono sistematicamente membri del loro gruppo, anche quando questi si comportano peggio di membri di altri gruppi.

In esperimenti controllati, gli scimpanzé aiutano più volentieri membri del loro gruppo che estranei, condividono più cibo con chi gli assomiglia, e mostrano più aggressività verso chi è diverso. Hanno perfino qualcosa che assomiglia al razzismo: discriminano sistematicamente contro scimpanzé di altre popolazioni, anche quando non c'è alcuna competizione per le risorse.

È lo stesso pattern che ritroviamo negli umani, solo amplificato. Noi non discriminiamo solo contro altre popolazioni: discriminiamo contro altre razze, altre nazioni, altre religioni, altri partiti politici, perfino altre squadre di calcio. Il meccanismo neurobiologico è identico, ma le sue manifestazioni sono diventate infinitamente più complesse e articolate.

La cosa inquietante è che questo pregiudizio per il gruppo di appartenenza si attiva automaticamente e inconsciamente. Gli scimpanzé non "decidono" di essere prevenuti: semplicemente lo sono, per programmazione evolutiva. E lo stesso vale per noi. Non scegliamo consciamente di preferire chi ci assomiglia: è il nostro cervello primitivo che fa quella scelta per noi, prima ancora che ce ne rendiamo conto.

\medskip

Ma c'è un aspetto ancora più sottile di questa eredità cognitiva condivisa: la gestione delle gerarchie sociali. Negli scimpanzé, come in tutte le società di primati, esistono gerarchie di dominanza molto rigide. Ogni individuo sa esattamente qual è il suo posto nella scala sociale, e si comporta di conseguenza.

Gli scimpanzé di rango inferiore si sottomettono spontaneamente a quelli di rango superiore, anche quando non c'è alcuna minaccia immediata. Mostrano segni di stress quando devono interagire con dominanti, e evitano il contatto visivo diretto. Al contrario, gli scimpanzé dominanti assumono posture di controllo, occupano più spazio, e hanno accesso prioritario alle risorse.

Tutto questo vi suona familiare? Dovrebbe. È esattamente quello che succede in ogni ufficio, ogni scuola, ogni organizzazione umana. Abbiamo gli stessi pattern di sottomissione e dominanza, gli stessi segnali non verbali, la stessa ansia gerarchica. La differenza è che noi abbiamo sviluppato sistemi di gerarchie molto più complessi e astratti.

Mentre gli scimpanzé competono per il cibo e i partner sessuali, noi compeitiamo per stipendi, titoli professionali, followers sui social media, marche di auto, quartieri residenziali. Ma il software neurobiologico che gestisce queste competizioni è identico. È lo stesso cervello primitivo che cerca di stabilire la propria posizione nella gerarchia sociale, solo che ora la gerarchia include miliardi di persone e infiniti parametri di valutazione.

\bigskip

Questo ci porta alla questione centrale: cosa succede quando un cervello primitivo ottiene poteri non primitivi?

Gli scimpanzé vivono in gruppi di 30-60 individui. Il loro mondo sociale è limitato, concreto, immediato. Le loro decisioni hanno conseguenze locali e temporanee. Se un capo scimpanzé prende una decisione stupida, nel peggiore dei casi può danneggiare il suo gruppo per una stagione.

Noi umani abbiamo lo stesso cervello sociale, ma viviamo in un mondo di 8 miliardi di persone interconnesse. Le nostre decisioni possono avere conseguenze globali e permanenti. Un presidente può scatenare una guerra mondiale, un CEO può causare una crisi economica, uno scienziato può creare un'arma di distruzione di massa, un influencer può diffondere teorie cospirative a milioni di persone.

È come se avessimo dato a uno scimpanzé i poteri di Superman. Il cervello è sempre quello: programmato per gestire problemi semplici in gruppi piccoli. Ma i poteri sono diventati planetari.

Il risultato è che spesso affrontiamo problemi del XXI secolo con intuizioni del Pleistocene. Reagiamo al cambiamento climatico come se fosse una minaccia tribale immediata (ovvero lo ignoriamo perché è troppo astratto). Gestiamo l'economia globale con gli stessi istinti che usavamo per dividere la selvaggina nella savana. Facciamo diplomazia internazionale con la mentalità di chi negoziava territori di caccia.

\medskip

Prendiamo un esempio concreto: la gestione delle crisi finanziarie. Quando scoppia una bolla speculativa, i mercati reagiscono esattamente come un gruppo di scimpanzé in preda al panico: tutti corrono nella stessa direzione, senza pensare se sia la direzione giusta.

È il bias del gregge applicato alla finanza globale. Il cervello primitivo valuta la situazione così: "Se tutti stanno vendendo, deve esserci un pericolo. Meglio vendere anch'io, per sicurezza." È la stessa logica che porta un gruppo di scimpanzé a fuggire tutti insieme quando uno di loro grida l'allarme, anche se l'allarme è falso.

La differenza è che quando gli scimpanzé corrono tutti nella stessa direzione per un falso allarme, al massimo sprecano energia e perdono tempo. Quando gli esseri umani "corrono tutti nella stessa direzione" sui mercati finanziari, possono distruggere l'economia di intere nazioni e mandare in povertà milioni di persone.

È lo stesso cervello, la stessa logica, lo stesso errore. Ma le conseguenze sono amplificate per la potenza della nostra tecnologia e l'interconnessione della nostra società.

\bigskip

Ma forse l'esempio più drammatico di questa discrepanza tra neurobiologia primitiva e potere moderno lo vediamo nella politica contemporanea. I leader politici sono esseri umani con cervelli di primati, ma hanno il potere di decidere le sorti di milioni di persone.

I bias cognitivi che in un capo scimpanzé potrebbero al massimo causare qualche conflitto locale, in un leader umano possono scatenare guerre, crisi economiche, disastri ambientali. L'overconfidence che spinge uno scimpanzé alfa a sfidare un rivale può spingere un presidente a invadere un paese. L'avversione alle perdite che fa sì che uno scimpanzé si aggrappi al territorio può portare un leader a prolungare una guerra disastrosa piuttosto che ammettere la sconfitta.

Il bias di conferma che porta uno scimpanzé a ignorare segnali di pericolo perché non si adattano alle sue aspettative può portare un leader politico a ignorare rapporti di intelligence cruciali. L'effetto ancoraggio che influenza le decisioni di uno scimpanzé su dove cercare cibo può influenzare le decisioni di un leader su come allocare il budget nazionale.

E il pregiudizio per il gruppo di appartenenza, che negli scimpanzé si limita a preferire membri del proprio clan, negli umani si trasforma in nazionalismo, razzismo, fanatismo religioso, e tutte le forme di "noi contro loro" che hanno causato la maggior parte delle sofferenze della storia umana.

\medskip

La cosa più inquietante di tutto questo è che non possiamo semplicemente "crescere" oltre questi bias. Non sono difetti di maturità o educazione che possiamo superare con più esperienza o migliore formazione. Sono parte dell'hardware neurobiologico che condividiamo con tutti i primati.

Un CEO di una multinazionale può avere tre lauree, vent'anni di esperienza, e consiglieri brillanti, ma quando deve prendere una decisione sotto pressione, il suo cervello reagisce esattamente come quello di uno scimpanzé che deve decidere se sfidare il capo branco. Gli stessi ormoni, gli stessi circuiti neurali, gli stessi bias cognitivi.

La differenza è che lo scimpanzé, nel peggiore dei casi, può perdere qualche banana. Il CEO può mandare in bancarotta un'azienda e licenziare migliaia di persone.

È per questo che le nostre società moderne sono piene di meccanismi di controllo, contrappesi, procedure burocratiche, e sistemi di supervisione che sembrano irrazionalmente complicati. Non sono complicati per cattiveria o inefficienza: sono complicati perché devono compensare il fatto che stiamo dando poteri planetari a cervelli progettati per gestire gruppi di 50 cacciatori-raccoglitori.

\bigskip

Ma c'è anche una lezione di umiltà in tutto questo. Quando critichiamo i leader politici per le loro decisioni irrazionali, quando ci lamentiamo delle follie dei mercati finanziari, quando ci stupiamo della stupidità delle folle, stiamo essenzialmente criticando la natura umana stessa.

Quei leader, quei trader, quelle folle sono fatti dello stesso materiale cognitivo di cui siamo fatti noi. Hanno gli stessi bias, gli stessi punti ciechi, le stesse limitazioni neurobiologiche. L'unica differenza è che le loro decisioni hanno conseguenze più visibili e drammatiche delle nostre.

Ma tutti noi, ogni giorno, prendiamo decisioni con lo stesso cervello di scimpanzé. Scegliamo lavori, partner, investimenti, stili di vita basandoci su intuizioni sviluppate per un mondo che non esiste più da decine di migliaia di anni. La differenza è che la maggior parte delle nostre decisioni sbagliate rimane privata e locale.

È un po' come criticare qualcuno per aver fatto un incidente stradale quando anche noi guidiamo la stessa macchina con gli stessi freni difettosi. La differenza è che lui stava guidando in autostrada a 200 all'ora, mentre noi andiamo piano nel quartiere. Ma la macchina è la stessa.

\bigskip

Allora, cosa possiamo fare con questa consapevolezza? La prima cosa è smettere di aspettarci che gli esseri umani - inclusi noi stessi - si comportino come computer perfettamente razionali. Non lo faremo mai, perché non siamo mai stati progettati per farlo.

La seconda è progettare sistemi, istituzioni, e processi decisionali che tengano conto dei nostri bias invece di ignorarli. Invece di aspettarci che i leader siano immuni all'overconfidence, possiamo creare meccanismi che li costringano a consultare pareri dissidenti. Invece di sperare che i mercati siano razionali, possiamo introdurre circuit breaker e regole che limitino i danni dei comportamenti irrazionali di massa.

La terza è sviluppare un po' di compassione per la condizione umana. Siamo tutti scimpanzé con pretese di grandezza, cervelli primitivi in un mondo moderno, animali sociali che cercano di gestire problemi per i quali non siamo stati progettati.

E forse, alla fine, questa è anche la cosa più bella dell'essere umani: siamo riusciti a creare una civiltà incredibilmente sofisticata nonostante - o forse proprio grazie - ai nostri limiti cognitivi. Siamo scimpanzé che sono andati sulla Luna, primati che hanno dipinto la Cappella Sistina, animali che si interrogano sui misteri dell'universo.

Non male per un cervello progettato per raccogliere bacche e evitare i predatori. Forse dovremmo essere più indulgenti con noi stessi, e più realistici su quello che possiamo aspettarci dalla nostra specie. Dopo tutto, per essere delle scimmie, non ce la stiamo cavando male.